% META: {"id": "1SPE-DERLOCAL-EV-A-corrige", "chapitre": "1SPE-DERIVATION-LOCAL", "type_objet": "corrige_evaluation", "capacites_codes": ["C1", "C2", "C3", "C4", "C5"], "status": "generated", "evaluation_ref": "1SPE-DERLOCAL-EV-A"}

\section*{Corrigé — Évaluation A — Dérivation locale}

%---------------------------------------------------------------
\subsection*{Exercice 1 — 6 points}
%---------------------------------------------------------------

\baremeIndicatif{Q1 : 1,5 pt ; Q2 : 2 pts ; Q3 : 1 pt ; Q4 : 1,5 pt}

\textbf{Question 1.}
$f(1) = 1 - 4 + 5 = 2$ et $f(4) = 16 - 16 + 5 = 5$.
\[
\tau = \frac{f(4)-f(1)}{4-1} = \frac{5-2}{3} = 1.
\]

\textbf{Question 2.}
$f(3) = 9 - 12 + 5 = 2$.
\begin{align*}
f(3+h) &= (3+h)^2 - 4(3+h) + 5 \\
        &= 9 + 6h + h^2 - 12 - 4h + 5 = h^2 + 2h + 2.
\end{align*}
\[
\frac{f(3+h)-f(3)}{h} = \frac{h^2+2h+2-2}{h} = \frac{h^2+2h}{h} = h+2.
\]

\textbf{Question 3.}
$f'(3) = \lim_{h \to 0}(h+2) = 2.$

\textbf{Question 4.}
$f'(3) = 2$ signifie que la tangente à la courbe de $f$ au point d'abscisse $3$ a pour pente $2$ : la courbe monte localement avec un coefficient directeur de $2$.

%---------------------------------------------------------------
\subsection*{Exercice 2 — 5 points}
%---------------------------------------------------------------

\baremeIndicatif{Q1 : 1 pt ; Q2 : 2 pts ; Q3 : 1 pt ; Q4 : 1 pt}

\textbf{Question 1.}
$g'(1) = 3 \times 1^2 = 3.$

\textbf{Question 2.}
$g(1) = 1$. La tangente $T_1$ en $x=1$ :
\[
T_1 : y = g'(1)(x-1) + g(1) = 3(x-1) + 1 = 3x - 2.
\]

\textbf{Question 3.}
$T_1(0) = 3 \times 0 - 2 = -2$. Le point $(0;-2)$ est bien sur $T_1$. \checkmark

\textbf{Question 4.}
Tangente horizontale $\Leftrightarrow g'(a) = 0 \Leftrightarrow 3a^2 = 0 \Leftrightarrow a = 0$.
La tangente est horizontale au point d'abscisse $0$.

%---------------------------------------------------------------
\subsection*{Exercice 3 — 5 points}
%---------------------------------------------------------------

\baremeIndicatif{Q1 : 1 pt ; Q2 : 1,5 pt ; Q3 : 1,5 pt ; Q4 : 1 pt}

\textbf{Question 1.}
$f(2) = 7$ (ordonnée du point $A$) et $f'(2) = -3$ (pente de la tangente en $A$).

\textbf{Question 2.}
$T : y = f'(2)(x-2)+f(2) = -3(x-2)+7 = -3x + 13.$

\textbf{Question 3.}
$f(2{,}04) \approx f(2) + f'(2) \times 0{,}04 = 7 + (-3) \times 0{,}04 = 7 - 0{,}12 = 6{,}88.$

\textbf{Question 4.}
C'est une valeur approchée (symbole $\approx$). L'approximation linéaire remplace la courbe par sa tangente au voisinage du point ; l'écart avec la valeur exacte est d'autant plus petit que $|h|$ est petit.

%---------------------------------------------------------------
\subsection*{Exercice 4 — 4 points}
%---------------------------------------------------------------

\baremeIndicatif{Q1 : 1,5 pt ; Q2 : 1,5 pt ; Q3 : 1 pt}

\textbf{Question 1.}
$h(2) = \dfrac{1}{2}$ et $h(3) = \dfrac{1}{3}$.
\[
\tau = \frac{h(3)-h(2)}{3-2} = \frac{\frac{1}{3}-\frac{1}{2}}{1} = \frac{-\frac{1}{6}}{1} = -\frac{1}{6} \approx -0{,}167.
\]

\textbf{Question 2.}
$h'(2) = -\dfrac{1}{4}$.
\[
h(2{,}05) \approx h(2) + h'(2) \times 0{,}05 = \frac{1}{2} + \left(-\frac{1}{4}\right) \times 0{,}05 = 0{,}5 - 0{,}0125 = 0{,}4875.
\]

\textbf{Question 3.}
Valeur exacte : $\dfrac{1}{2{,}05} = \dfrac{100}{205} = \dfrac{20}{41} \approx 0{,}48780\ldots$

Approximation : $0{,}4875$. L'erreur est de l'ordre de $0{,}0003$, soit environ $0{,}06\,\%$.

% BEGIN-VERIFY
% from sympy import symbols, diff, simplify, solve, Rational, limit
% x, h = symbols('x h')
% # ---- Exercice 1 : f(x) = x^2 - 4x + 5 ----
% f = x**2 - 4*x + 5
% assert f.subs(x, 1) == 2 and f.subs(x, 4) == 5
% assert Rational(f.subs(x, 4) - f.subs(x, 1), 4 - 1) == 1
% assert f.subs(x, 3) == 2
% assert simplify(f.subs(x, 3 + h) - (h**2 + 2*h + 2)) == 0
% taux = simplify((f.subs(x, 3 + h) - f.subs(x, 3))/h)
% assert simplify(taux - (h + 2)) == 0
% assert limit(taux, h, 0) == 2
% assert diff(f, x).subs(x, 3) == 2
% # ---- Exercice 2 : g(x) = x^3 ----
% g = x**3
% gp = diff(g, x)
% assert gp.subs(x, 1) == 3 and g.subs(x, 1) == 1
% T1 = gp.subs(x, 1)*(x - 1) + g.subs(x, 1)
% assert simplify(T1 - (3*x - 2)) == 0
% assert T1.subs(x, 0) == -2
% assert solve(gp, x) == [0]
% # ---- Exercice 3 : lecture graphique f(2)=7, f'(2)=-3 ----
% f2, fp2 = 7, -3
% T = fp2*(x - 2) + f2
% assert simplify(T - (-3*x + 13)) == 0
% assert f2 + fp2*Rational(4, 100) == Rational(688, 100)
% # ---- Exercice 4 : h(x) = 1/x ----
% hh = 1/x
% assert hh.subs(x, 2) == Rational(1, 2) and hh.subs(x, 3) == Rational(1, 3)
% assert simplify((hh.subs(x, 3) - hh.subs(x, 2))/(3 - 2) + Rational(1, 6)) == 0
% hp = diff(hh, x)
% assert hp.subs(x, 2) == Rational(-1, 4)
% assert Rational(1, 2) + Rational(-1, 4)*Rational(5, 100) == Rational(4875, 10000)
% assert simplify(hh.subs(x, Rational(205, 100)) - Rational(20, 41)) == 0
% END-VERIFY
